\begin{textsample}{POS Dim 3 – ms – Score 13.00 – ms/ms\_inf\_e\_ef\_29.txt}  \label{ex:f3_pos_233}
7.4.4 \textbf{Escuta}, Fala, Pensamento e Imaginação As DCNEI ( 2009 ) propõem que sejam garantidas nas propostas pedagógicas experiências que promovam a imersão das \textbf{crianças} nas diferentes linguagens e o progressivo domínio por elas de vários gêneros e formas de expressão : gestual, verbal, \textbf{plástica}, \textbf{dramática} e musical.
As \textbf{crianças} vivem de forma intensa suas experiências, por meio das interações, \textbf{brincadeiras}, gestualidades, movimentos, linguagem verbal ou da Língua Brasileira de Sinais ( LIBRAS ), significam e dão novos significados ao mundo à sua moda : correm e pulam ; \textbf{conversam} e recontam histórias ; leem e escrevem ; \textbf{contam}, cantam e dançam ; pintam e desenham ; observam e experimentam ; imitam, imaginam e \textbf{brincam} ; choram e riem ; brigam e compartilham.
Desse modo : > As manifestações linguageiras das \textbf{crianças} e dos artistas convidam a reorganizar o mundo e experimentá -lo em outras versões, mediados pelos corpos que se mexem, que nem sempre falam com palavras e letras, mas que tanto dizem, provocando a conhecer o desconhecido ao mesmo tempo em que se constroem outros lugares de experiências, estranhando e conhecendo a todo instante ( GOBBI, 2010, p. 2 ).
Por isso, as \textbf{crianças} necessitam das diferentes linguagens para atuarem socialmente e manifestarem seus pensamentos, imaginação, ideias, sentimentos, hipóteses, \textbf{desejos}.
Essas manifestações transformadas em experiências potencializam a apropriação do conhecimento, a organização do pensamento, a capacidade criativa, expressiva e comunicativa e a participação na cultura.
Assim, a Educação \textbf{Infantil} é um lugar, por excelência, de acesso às diferentes linguagens : a musical, a \textbf{plástica}, a corporal e a pictórica, as quais são elementos de mediação, interlocução e interação entre as \textbf{crianças} e os \textbf{adultos} e devem ser utilizadas a favor da expressão, imaginação e criação das \textbf{crianças}.
Desse modo, a articulação das diferentes linguagens promove -lhes a aprendizagem e o desenvolvimento.
As diferentes linguagens articuladas marcam uma rede de interações que integra a \textbf{criança} ao seu meio sociocultural e histórico.
Por meio da fala, do pensamento, da \textbf{escuta} e da imaginação, ela apreende a realidade, interage nela e com ela, transforma -a e, consequentemente, desenvolve -se, transforma -se, elabora conhecimentos e amplia seus referenciais de mundo e cultura. 7.4.5 Espaços, Tempos, Quantidades, Relações e Transformações O trabalho intencional das experiências mediadas pelo professor abre um leque de possibilidades para que as \textbf{crianças} se apropriem dos conhecimentos matemáticos e os relativos ao mundo social e natural, os quais permitem novas formas de pensar para o desenvolvimento da cultura científica.
A valorização dos conhecimentos prévios deve ser uma premissa de trabalho, porém Arce alerta : > A chave para pensarmos o trabalho nas salas de educação \textbf{infantil} reside na vida cotidiana das \textbf{crianças}, quando o docente consegue integrar os conteúdos científicos aos conhecimentos presentes no dia a dia delas.
Quando esses conhecimentos ganham nova forma é que a escola contribui qualitativamente para desenvolvimento intelectual das \textbf{crianças} ( ARCE, 2018, p.110 ).
As \textbf{crianças} são curiosas, têm interesse pelo desconhecido ; ao propor a familiarização com os conhecimentos presentes nesse campo, o professor deve evitar o reforço de estereótipos culturais e possibilitar experiências específicas da região em que vivem e ampliar o acesso a diferentes práticas sociais.
Importante destacar que é preciso atenção aos encaminhamentos junto às \textbf{crianças}, pois pelo receio de antecipar conclusões, o professor pode negar informações, cuja compreensão se dará a partir do entendimento da capacidade mental que cada uma possui naquele momento.
Experiências relacionadas aos saberes e conhecimentos que incentivem a postura investigativa e agucem a \textbf{curiosidade} e capacidade de questionar, conhecer, argumentar, levantar e checar hipóteses certamente contribuirão para o desenvolvimento integral e intelectual das \textbf{crianças}.

% matched lemmas: adulto, brincadeira, brincar, contar, conversar, criança, curiosidade, desejo, dramático, escuta, infantil, plástico
\end{textsample}
